\appendix
\section{Supplementary Experimental Details}
\label{app:configs}

The following sections provide supplementary experimental details, exact hyperparameter configurations, and verification protocols utilized across the main body of this paper. These technical details are included to ensure full reproducibility of the computational experiments without obstructing the primary narrative.

\subsection{Per-Experiment Training Budgets}

To systematically map the PINN failure landscape, we executed highly varied computational budgets dependent upon the specific demands of each experimental protocol. The default budget baseline consisted of $40,000$ optimization epochs, but was heavily modified for dense grid searches, variance decompositions, and large statistical ensemble aggregations.

\begin{table}[htpb]
    \centering
    \caption{Training budget and collocation density configurations for representative primary experiments.}
    \label{tab:app_budgets}
    \begin{tabular}{@{}lllllrll@{}}
        \toprule
        Experiment & PDE & $N_{epochs}$ & $N_{col}$ & $N_{IC}$ & $N_{BC}$ & $N_{seeds}$ & Device \\
        \midrule
        Exp 1 (Spectral) & Advection & $40,000$ & $1000$ & $100$ & $100$ & $10$ & NVIDIA RTX 3050 6GB \\
        Exp 5 (Hessian)  & Advection & $20,000$ & $2500$ & $250$ & $250$ & $5$ & NVIDIA RTX 3050 6GB \\
        Exp 7 (Sampling) & Helmholtz & $10,000$ & $2000$ & $-$ & $400$ & $20$ & NVIDIA RTX 3050 6GB \\
        Exp 8 (Starvation) & Helmholtz & $50,000$ & Sweep & $150$ & $150$ & $3$ & NVIDIA RTX 3050 6GB \\
        Exp 9 (BC Ratio) & Helmholtz & $10,000$ & Sweep & $100$ & Sweep & $3$ & NVIDIA RTX 3050 6GB \\
        Exp 11 (Multi-Scale) & Allen-Cahn & $30,000$ & $1000$ & $100$ & $100$ & $3$ & NVIDIA RTX 3050 6GB \\
        Exp 13 (Init) & Advection & $40,000$ & $1000$ & $100$ & $100$ & $50$ & NVIDIA RTX 3050 6GB \\
        Exp 14 (Optimizer) & Advection & $40,000$ & $1500$ & $200$ & $200$ & $3$ & NVIDIA RTX 3050 6GB \\
        Exp 25 (Phase)   & All PDEs & $20,000$ & Grid & Grid & Grid & $1$ & NVIDIA RTX 3050 6GB \\
        Exp 27 (ANOVA)   & Adv/Burg & $8{,}000$ & $3000$ & $150$ & $150$ & $1{,}134$\,\textsuperscript{$\dagger$} & NVIDIA RTX 3050 6GB \\
        \bottomrule
    \end{tabular}
\end{table}
\noindent\textsuperscript{$\dagger$}Total runs across the full factorial grid: $648$ (Advection: $4\beta \times 3 \text{arch} \times 3\text{lr} \times 3W \times 3\text{init} \times 2\text{seeds}$) $+ 486$ (Burgers: $3\nu \times 3 \times 3 \times 3 \times 3 \times 2$) $= 1{,}134$.

\subsection{Reference Solution Methods}
\label{app:reference}

Because standard PINN evaluation inherently requires calculation of the relative $\Ltwo$ error against a ground-truth physical solution, highly accurate reference solvers were required for all tested PDEs. 

For the 1D Advection equation ($u_t + \betaAdv u_x = 0$), the analytical solution $u_{\text{exact}}(x,t) = \sin(x - \betaAdv t)$ was explicitly utilized, sampled across a highly dense, uniform $1024 \times 1024$ spatio-temporal evaluation grid for exact error computation.

For the Viscous Burgers equation ($u_t + u u_x = \nuBurg u_{xx}$), no exact general analytical solution was universally suitable across all evaluated stiff regimes. We therefore generated high-fidelity ground truth numerical solutions utilizing a standard Crank-Nicolson finite difference scheme. To properly resolve the severe shock fronts at low viscosities, we enforced a high spatial resolution of $dx = 1/2048$ combined with a dynamic, adaptive time-stepping protocol strictly defined by $dt = \min(\text{CFL} \cdot dx/|u|_{\text{max}}, 0.4 \cdot dx^2/(2\nuBurg))$. Results for Burgers at $\nuBurg < 0.001$ are excluded from the main analysis because the reference solver stability degrades at this regime.

For the 2D Helmholtz problem, we utilized the mathematically exact method of manufactured solutions. The localized forcing function $f(x,y)$ was structurally chosen so that the target analytical function $u = \sin(\pi x)\sin(\pi y)$ becomes the exact, zero-error solution, allowing direct computation of the network error independently of numerical integration bounds.

\subsection{Exp 5: Multi-Seed Hessian Protocol}
\label{app:hessian}

To investigate the local geometry of the optimization landscape in Exp.~5 (The Sharp-Trap Paradox), we required precise estimation of the dominant Hessian eigenvalue ($\lambda_{\text{max}}$). We employed an identical computational protocol across $N_{\text{HESSIAN\_RUNS}} = 5$ distinct, isolated random seeds. The dominant eigenvalue was computed efficiently via power iteration utilizing $25$ steps, mathematically avoiding the requirement to explicitly form the prohibitively large, memory-intensive $P \times P$ Hessian matrix.

Failed models were systematically generated using an aggressive Normal initialization ($\sigma=1.0$), with a genuine failure criterion strictly defined as yielding a relative $\Ltwo > 0.1$. The successful models yielded a mean $\lambda_{\text{max}}$ of $3362.4 \pm 810.9$. The mean Hessian maximum eigenvalue of the failed model ($93968923.2 \pm 97965534.6$, $n=5$ runs) exceeded that of the successful model ($3362.4 \pm 810.9$) by a factor of $27947$; however, the $2\sigma$ confidence intervals overlap ($-101962146.0$ to $289899992.4$ versus $1740.6$ to $4984.2$), indicating the difference does not reach the pre-specified significance threshold at $n=5$. The directional finding — failed models occupying sharper landscape regions — is consistent across all 5 runs, but the high relative variance of the Hessian estimator (relative SD = $104.3\%$) precludes a strong statistical claim at this sample size. Larger $n$ is required for formal confirmation.

\subsection{Exp 14: Budget Equalization Protocol}
\label{app:budget}

Directly comparing the empirical efficiency of first-order and second-order optimizers (e.g., Adam vs.\ L-BFGS) requires a rigorously equalized computational budget, as measuring uncalibrated "epochs" across distinct algorithms is inherently misleading. In Exp.~14, we systematically constrained the total number of physical gradient evaluations.

For the standard Adam optimizer baseline, training inherently consisted of $40,000$ sequential epochs, equivalent to exactly $40,000$ forward-backward gradient steps. To strictly equalize this physical budget for L-BFGS, the solver was mathematically constrained to receive exactly $N_{\text{LBFGS\_STEPS}} = N_{\text{EPOCHS}} / \text{LBFGS\_MAX\_ITER}$ steps. Because we explicitly configured $\text{LBFGS\_MAX\_ITER} = 20$, the L-BFGS optimizer cleanly executed $40,000 / 20 = 2000$ outer steps, each containing up to $20$ internal closure evaluations. Consequently, $2000 \text{ steps} \times 20 \text{ closure calls} = 40,000$ total gradient evaluations, precisely mirroring the exact computational cost of the Adam baseline.

\subsection{Exp 27: ANOVA Design Specification}
\label{app:anova}

The comprehensive variance decomposition executed in Exp.~27 utilized a robust factorial design to isolate structural influences versus hyperparameter sensitivities. The explicitly tested factor levels are detailed comprehensively below.

\begin{table}[htpb]
    \centering
    \caption{Explicit factor levels evaluated in the $1{,}134$-run factorial ANOVA design ($567$ unique cells $\times$ $2$ seeds). The stiffness factor is PDE-specific: $S_1$ ($\betaAdv$) for Advection and $S_2$ ($\nuBurg$) for Burgers; the architecture factor $S_3$ is shared across both PDEs.}
    \label{tab:app_anova}
    \begin{tabular}{@{}lll@{}}
        \toprule
        Group & Factor & Discrete Levels \\
        \midrule
        Structural & Stiffness ($\betaAdv$ or $\nuBurg$) & Advection: $[1, 10, 30, 50]$ | Burgers: $[0.1, 0.01, 0.001]$ \\
        Structural & Architecture Class & Deep MLP (Tanh), Deep MLP (SiLU), Shallow MLP (Tanh) \\
        Hyperparameter & Learning Rate & $10^{-4}$, $10^{-3}$, $5\times10^{-3}$ \\
        Hyperparameter & Network Width & $32$, $64$, $128$ neurons/layer \\
        Hyperparameter & Initialization & Glorot Uniform, He Normal, Orthogonal \\
        \bottomrule
    \end{tabular}
\end{table}

For maximum analytical rigor, the effect size for each individual factor was explicitly calculated using the standard unbiased $\etasq$ estimator formula: $\etasq = SS_{\text{factor}} / SS_{\text{total}}$, where $SS_{\text{factor}}$ represents the sum of squares strictly attributable to the given parameter variance factor, and $SS_{\text{total}}$ represents the total aggregated outcome variance across the massive experimental grid. 

The additive SS decomposition used here does not account for factor interactions. A full factorial ANOVA with interaction terms would require a substantially larger design and is left for future work.

\subsection{Supplementary Material: Exact Seed Values}
\label{app:seeds}

To facilitate bit-for-bit reproduction of the $n=1$ and small-$n$ experiments, we explicitly record the exact random seeds utilized for initialization and dataset generation:

\begin{itemize}
    \item \textbf{Exp 1 (Spectral Bias)}: Seed 42
    \item \textbf{Exp 2 (Activation Study)}: Seed 42
    \item \textbf{Exp 3 (NTK Analysis)}: Seed 42
    \item \textbf{Exp 4 (Gradient Pathology)}: Seed 101
    \item \textbf{Exp 6 (Gradient Flow)}: Seed 101
    \item \textbf{Exp 10 (Stiffness Failure)}: Seed 42
    \item \textbf{Exp 16 (Causality)}: Seed 42
    \item \textbf{Exp 25 (Phase Space)}: Seed 42
    \item \textbf{Exp 26 (Conservation Laws)}: Seed 42
\end{itemize}

Where $n=3$ or $n=5$ aggregates are reported (e.g., Exp 8, Exp 5), the sequential seeds $\{42, 43, 44, 45, 46\}$ were systematically applied. For the massively replicated distributions ($n=1{,}134$, $n=400$), random states were initialized uniformly from a master seed of $0$ utilizing NumPy's \texttt{RandomState} generator.
